\section{Introduction}
\label{sec:intro}

% P1: Hardware-aware software optimization and agent-generated kernels.
Efficient numerical software depends on kernels that exploit the execution and memory organization of the target processor. Developing these kernels requires choices about data layout, vectorization, blocking, and instruction scheduling that are difficult to make uniformly across hardware platforms. Recent work investigates language-model agents as a means of exploring this implementation space. KernelBench evaluates generated GPU kernels against reference tensor computations, while Autocomp combines hardware knowledge with correctness and performance feedback to search for faster accelerator programs~\cite{kernelbench,autocomp}. These systems demonstrate opportunities for automating kernel development, while making the reliability of generated implementations an increasingly practical concern.

% P2: CPU SIMD, architectural differences, and the intrinsic interface.
We focus on SIMD and vector kernels written using C intrinsics. Vector extensions in modern CPUs expose different execution models: Arm NEON operates on fixed-width vectors, whereas RISC-V Vector (RVV) provides runtime selection of the number of active elements~\cite{acle,rvv}. These differences affect loop structure, tail processing, and memory accesses. Architectures also differ in their numerical operations, even when the operations have similar names. RVV has attracted dedicated work on porting existing intrinsic code, including SALTyRN and IntrinTrans~\cite{saltyrn,intrintrans}. Compiler intrinsic interfaces make such work possible within C: programmers combine ordinary control flow and pointer manipulation with architecture-specific operations that the compiler lowers to target instructions. Intrinsics expose hardware capabilities without requiring an entire kernel to be written in assembly, but do not make its behavior portable by construction.

% P3: Existing validation must be described accurately.
This interface also gives agents a concrete programming environment in which to apply instruction descriptions and optimization guidance. SALTyRN translates NEON kernels to RVV and applies an Autocomp-based optimization stage~\cite{saltyrn}. Correctness remains a separate obligation. KernelBench's maintainers document generated programs that exploit evaluation weaknesses, including output-memory reuse and misleading timing~\cite{kernelbench-eval}. Existing systems address correctness in several ways: Autocomp tests candidates against reference implementations, IntrinTrans uses compile-and-test feedback, SALTyRN performs bounded symbolic equivalence checking, and LLM-Vectorizer uses Alive2 to validate vectorization~\cite{autocomp,intrintrans,saltyrn,llm-vectorizer}. Our question is how to establish a whole-kernel guarantee for the stated input domain, including loop iterations and memory accesses, while retaining the implementation produced by optimization. Such a guarantee can reject functional deviations that escape tests; it does not by itself prevent timing manipulation or other attacks on a benchmark harness.

% P4: Three abstraction levels, with concrete prior work.
The semantic interface for this verification can be placed at several levels. At the machine-code level, Sail supplies detailed ISA definitions, and Islaris supports program proofs against authoritative instruction semantics~\cite{sail,islaris}. At the intermediate-representation level, Alive2 checks LLVM transformations, while Rake synthesizes semantically equivalent vector instruction sequences for Hexagon HVX from Halide expressions~\cite{alive2,rake}. An IR-based comparison must account for the lowering of intrinsics and the memory semantics of the selected IR. C intrinsics offer another boundary: an executable C interpretation of each supported intrinsic can expose its behavior to a C program verifier, allowing the surrounding loops and pointer operations to remain in the language in which the kernel was written.

% P5: Scope-based design choice, not a claim that ISA proofs are impossible.
We choose this C-level boundary to study \emph{functional equivalence under explicit calling contracts}. The observations are the specified outputs and memory effects of ordinary kernel calls, under a stated floating-point and vector configuration. Interrupt handling, privileged execution, and timing are outside this observation. NaNs, infinities, integer overflow conditions, and invalid memory accesses still require treatment whenever they affect an admitted call. A binary-level proof can support a stronger connection to machine execution, but also requires relating machine states and establishing the relevant execution environment. Our C-level design keeps the proof interface close to the intrinsic source and permits reuse of C verification infrastructure. It introduces its own obligations: the intrinsic interpretations and the preparation of verification inputs must faithfully represent the source. We therefore retain these correspondences as explicit trust boundaries.

% P6: Framework and present implementation boundaries.
We present \tool, a framework for checking semantic equivalence between intrinsic kernels through reusable C semantic shims and modular proofs. The design covers Arm NEON, RVV, x86 AVX, and Hexagon HVX through a common contract interface. \todo{Establish and report the supported intrinsic subset and admitted kernel coverage for each ISA; four complete backends are not yet an experimental result.} For each kernel pair, \tool associates intrinsic calls with their semantic definitions and expresses the required output relation in a public specification. It separates the implementation, specification, and proof annotations so that proof development need not rewrite the executable algorithm. QCP symbolically executes the annotated C programs and generates verification conditions, which are discharged through its automation and Rocq/Coq proofs~\cite{qcp}. The implementation proofs establish the contracts used in the final equivalence argument. An unresolved proof obligation remains unresolved unless a checked proof or a reproducible counterexample is obtained.

% P7: Human review and justified preconditions.
The workflow concentrates per-kernel review on the calling contract, verification evidence, and any proposed implementation changes. A compact public specification provides an entry point to this review; its imported mathematical definitions and memory predicates remain part of what must be justified. When a counterexample violates a documented assumption of the original caller, the specification may be strengthened to express that assumption, provided the revised domain still includes all legitimate calls. When the counterexample is feasible, the target implementation requires repair or the claimed equivalence must be rejected. Repairs are reviewed as executable changes and measured for their performance effect before the revised candidate is accepted. Shared semantic models and the verification toolchain require separate review, so this workflow does not presume that reading a short specification is sufficient to establish trust in the entire system.

% P8: Three applications of one relation, not three separate frameworks.
The same equivalence interface serves three uses. First, it validates optimizations within an ISA or translations across ISAs against a reference implementation, without making the optimizer itself part of the trusted reasoning. Second, it supports porting kernels from existing NEON or AVX implementations to target libraries whose coverage of the selected operations is incomplete. Here RVV and Hexagon HVX are target backends; the motivation is missing implementations in a particular library, rather than a claim that both architectures are new. Third, it supports semantic cross-checking of independently developed implementations of the same operation. Differential testing compares executions on selected inputs~\cite{mckeeman}; program equivalence checking, exemplified by SymDiff~\cite{symdiff}, makes the relationship itself the verification target. A concrete disagreement can then be classified as an implementation defect, an allowed numerical difference, an unsupported calling assumption, or a modeling error. This classification identifies whether maintenance requires a code repair, a clearer contract, or a correction to the verification model.

% Alternative short P8, retained as requested; not rendered.
% The framework supports validation of kernel optimizations, verified translation
% across ISAs, and cross-checking of existing implementations. The planned main
% study uses SALTyRN's 35 published final optimized RVV kernels. Additional studies
% target AVX-to-RVV/HVX ports, XNNPACK implementation comparisons, and a composed
% SIMD program. All proof counts, repairs, cost measurements, and speedups remain
% \todo{pending evaluation}; none is asserted as a completed contribution.

% P9: Evaluation subjects are not completed results.
Our evaluation is organized around these uses. The primary study targets the 35 final optimized RVV kernels published by SALTyRN, checking each against its corresponding NEON source rather than reusing results for earlier intermediate versions~\cite{final-kernels}. Further studies are planned for five AVX operations ported to RVV and HVX, cross-checking of corresponding XNNPACK implementations, and a small composed SIMD program, with softmax as a candidate. \todo{Report verified pairs out of 35, unchanged-source coverage, justified repairs, verification time and token cost, and speed relative to a fixed handwritten RVV baseline.} \todo{Report completed new ports, classified discrepancies with upstream status, and the selected program's composed proof and performance.} These measurements are needed to assess whether the framework preserves optimized code, extends across ISA interfaces, and provides useful diagnostic evidence at an acceptable verification cost.
